{\rednote{TODO: This chapter needs more structure and contains lots of results, which for which we just didn't find a better place yet.}}

\subsection{Basics on Bundles}

The following lemmas and parts of their proofs are direct adaptions from \cite[Chapter II.5]{Hartshorne}.

\begin{lemma}
  \label{affine-extension-of-sections}
  Let $X$ be an affine scheme and $f:\Gamma(X,\mc O_X)$.
  Let furthermore $\mathcal F : X \to \Mod R$ be a weakly quasi-coherent bundle of $R$-modules.
  \begin{itemize}
    \item[i)] Given $s:\Gamma (X,\mathcal F)$ with $s|_{D(f)} = 0$  there merely exists an $n>0$ such that $f^ns=0$.
    \item[ii)] For any section $t:\Gamma (D(f),\mathcal F)$ there merely exists an $n>0$
      such that $f^nt$ extends to a global section of $\mathcal F$.
  \end{itemize}
\end{lemma}

\begin{proof}
  Since $\mathcal F$ is weakly quasi-coherent we have $\Gamma (X,\mathcal F)_f = \Gamma (D(f),\mathcal F)$.
  This implies both the first and second part.
\end{proof}


\begin{lemma}
  \label{extension-of-sections-by-twisting}
  Let $X$ be a scheme and $\mathcal L$ a line bundle on $X$ with a given section $f:\Gamma (X,\mathcal L)$.
  Let $\mathcal F:X\to \Mod{R}_{\mathrm{wqc}}$ be a weakly quasi-coherent bundle of $R$-modules.
  Denote by $D(f) \subseteq X$ the open locus where $f$ does not vanish.
  \begin{itemize}
    \item[i)] Given $s:\Gamma (X,\mathcal F)$ with $s|_{D(f)} = 0$ there merely exists an $n>0$ such that $f^ns=0$,
      when considered as a global section of $\mathcal F \otimes \mathcal L^{\otimes n}$.
    \item[ii)] For any section $t:\Gamma (D(f),\mathcal F)$ there merely exists an $n>0$
      such that $f^nt$ extends to a global section of $\mathcal F \otimes \mathcal L^{\otimes n}$.
  \end{itemize}
\end{lemma}

This generalizes the above lemma by simultaneously allowing non-affine schemes, as well as twists.

\begin{proof}
  For the first part consider any affine open subscheme $U=\Spec A$ of $X$ that trivializes $\mathcal L$,
  i.e. suppose there is an equivalence  $\varphi:\mathcal L_U \to \mathcal O_U$ and set $g :\equiv \varphi(f|_U)$.
  Thus we have $D(f)\cap U = D(g)$ and $s|_{D(g)}=0$.
  By virtue of \Cref{affine-extension-of-sections} there is a $n>0$ such that $g^ns=0: \Gamma (U,\mathcal F)$.
  Using the equivalence $\mathcal F|_U\otimes\mathcal L^{\otimes n}|_U \to \mathcal F|_U$ this implies $f^ns=0$,
  when considered as a section of $\mathcal F\otimes\mathcal L^{\otimes n}$ over $U$.
  Since $X$ can be covered by finitely many such $U$, we finally obtain a natural number $n>0$ such that $f^ns=0$ as a global section of $\mathcal F\otimes\mathcal L^{\otimes n}$.

  {\color{red} TODO: write up part (ii)}
\end{proof}

\rednote{The following Lemma is also in \href{https://felix-cherubini.de/adequate.pdf}{this} draft.}

\begin{lemma}
  \label{tensor-gamma-commute-on-affines}
  Let $X$ be an affine scheme and suppose that $\mathcal F, \mathcal G : X \to \Mod R$ are pointwise weakly quasi-coherent bundles.
  Then the map
  \[
    \vartheta_X : \Gamma(X,\mathcal F) \otimes_{R^X}\Gamma(X,\mathcal G) \to \Gamma(X,\mathcal F \otimes \mathcal G), \;\;
    f \otimes g \mapsto (x \mapsto f(x) \otimes g(x))
  \]
  is an $R^X$-linear isomorphism.
\end{lemma}

For the proof we need the following vanishing criterion for tensors.

\begin{remark}
  \label{null-tensor-criterion}
  Let $M$ and $N$ be $A$-modules over an arbitrary commutative ring $A$.
  For a tensor $k = \sum_{j=1}^n x_j \otimes y_j : M \otimes_A N$ we have $k = 0 : M \otimes_A N$ if and only if there merely are $p \geq n$, $q \geq 0$ together with elements $x_{n+1},\hdots,x_p : M$, $z_1,\hdots,z_q : N$ and a matrix $G : A^{p \times q}$, such that with
  \begin{align*}
    S &\equiv (x_1 \ \cdots \ x_p) : M^{1 \times p}, \\
    T &\equiv (y_1 \ \cdots \ y_n \; 0 \ \cdots \ 0)^T : N^{p \times 1}, \\
    U &\equiv (z_1 \ \cdots \ z_q)^T : N^{q \times 1}
  \end{align*}
  we have $SG = 0 : M^{1 \times q}$ and $GU = T : N^{p \times 1}$.
\end{remark}

For finitely generated $M$, this is the null-tensor lemma \cite[Lemma IV-4.14]{lombardi-quitte}.
The arbitrary case follows from the criterion that $k=0 : M\otimes_A N$ if and only if, merely, there is a finitely generated submodule $M'\subseteq M$ containing $x_1,\ldots,x_n$ such that $k = 0 : M' \otimes_A N$; see also \cite[Exercise IV-4]{lombardi-quitte}.

\begin{proof}
  Say $X = \Spec A$.
  We start by showing injectivity of $\vartheta_X$.
  Let us first explain how this is a local problem.
  From commutative algebra we know that it is enough to find $a_1,\hdots,a_m:A$ with $(a_1,\hdots,a_m) = A$ such that $\vartheta_X[a_i^{-1}]$ is injective for all $i=1,\hdots,m$.
  As $\mc F$ and $\mc G$ are weakly quasi-coherent, we use \cite[Lemma 7.1.8]{draft} to see that the left vertical map in the commutative square
  \[ \begin{tikzcd}[column sep=large]
    \Gamma(X,\mc F)[a_i^{-1}] \otimes_{A[a_i^{-1}]} \Gamma(X,\mc G)[a_i^{-1}] \ar[d, "\simeq"]\ar[r, "\vartheta_X{[a_i^{-1}]}"]
    &
    \Gamma(X, \mc F \otimes \mc G)[a_i^{-1}] \ar[d]
    \\
    \Gamma(D(a_i),\mc F) \otimes_{R^{D(a_i)}}\Gamma(D(a_i),\mc G) \ar[r, "\vartheta_{D(a_i)}"]
    &
    \Gamma(D(a_i),\mc F \otimes \mc G)
  \end{tikzcd}
  \]
  is an isomorphism.
  Hence the injectivity of all $\vartheta_{D(a_i)}$ implies the injectivity of $\vartheta_X$.
  Together with the fact that $(a_1,\hdots,a_m) = A$ if and only if $X = D(a_1) \cup \hdots \cup D(a_m)$, this means in particular that, for a given element of the kernel of $\vartheta_X$, it is enough to prove its vanishing after replacing $X$ by the members of a finite affine cover. We shall use this when invoking Zariski-local choice below.

  Now let $k:\Gamma(X,\mc F) \otimes_A \Gamma(X,\mc G)$ be such an element, i.e.\ $\vartheta_X(k)=0$.
  Merely, we may find $f_1,\hdots,f_n : \Gamma(X,\mc F)$ and $g_1,\hdots,g_n : \Gamma(X,\mc G)$ such that $k = \sum_{j=1}^n f_j \otimes_A g_j$ and $\sum_{j=1}^n f_j(x) \otimes g_j(x) = 0  : \mc F_x \otimes \mc G_x$ for all $x:X$.
  By virtue of \Cref{null-tensor-criterion}, we merely get the datum
  \[
    p_x \geq n, \;\;
    q_x \geq 0, \;\;
    S_x : \mc F_x^{1 \times p_x}, \;\;
    T_x : \mc G_x^{p_x \times 1}, \;\;
    U_x : \mc G_x^{q_x \times 1}, \;\;
    G_x : R^{p_x \times q_x},
  \]
  such that $S_xG_x = 0$ and $G_xU_x = T_x$. Here, $S_x$, $T_x$ and $U_x$ are defined as in the remark.
  Explicitly, the first $n$ entries of $S_x$ are $f_1(x), \hdots, f_n(x)$  and the remaining are supplied datum by the remark.
  The first $n$ entries of $T_x$ are $g_1(x), \hdots, g_n(x)$ while the trailing entries are zero.
  Using Zariski-local choice while shrinking $X$ if necessary as explained in the first paragraph, we get functions $p,q : X \to \mathbb N$ such that $p_x \geq n$ for all $x:X$, together with dependent functions $S : (x:X) \to \mc F_x^{1 \times p_x}$, $T : (x:X) \to \mc G_x^{p_x \times 1}$, $U : (x:X) \to \mc G_x^{q_x \times 1}$ and $G : (x:X) \to R^{p_x \times q_x}$ such that $SG = 0$ and $GU = T$.
  The functions $p$ and $q$ are bounded by \cite[Theorem 3.3.3]{draft}.
  Furthermore, because equality in $\mathbb N$ is decidable and hence open, all fibers of $p$ and $q$ are open subschemes of $X$.
  It follows that $(\fib_p(n))_{n:\mathbb N}$ and $(\fib_q(n))_{n:\mathbb N}$ are both finite open covers of $X$.
  Again shrinking $X$, we thus may assume that $p$ and $q$ are constant.
  Therefore we have $S : \Gamma(X,\mc F)^{1 \times p}$, $T : \Gamma(X,\mc G)^{p \times 1}$, $U : \Gamma(X,\mc G)^{q \times 1}$ and $G : \Gamma(X, R)^{p \times q} \simeq A^{p \times q}$.

  By construction, the first $n$ entries of $S$ and $T$ are given by $f_1, \hdots, f_n$ and $g_1, \hdots, g_n$ respectively and the trailing entries of $T$ are zero.
  Therefore the reverse implication of \Cref{null-tensor-criterion} gives that $k = \sum_{j=1}^n f_j \otimes g_j = 0$\footnote{
    Alternatively, it is perhaps enlightening to do this calculation by foot:
    \[
      \textstyle
      \sum_{j=1}^n f_j \otimes g_j = \sum_{i=1}^p S_i \otimes T_i = \sum_{i,k} S_i \otimes G_{ik}U_k = \sum_k(\sum_i S_iG_{ik}) \otimes U_k = 0 \rlap{.}
    \]
    This is essentially the proof of the reverse implication of the remark.
  }.

  Lastly we verify surjectivity.
  Let $s:\Gamma(X,\mc F \otimes \mc G)$ and set $H :\equiv \Gamma(X,\mc F) \otimes_A \Gamma(X,\mc G)$.
  We know that $s(x) : \mc F_x \otimes \mc G_x$ merely is a finite sum of elementary tensors.
  Applying Zariski-local choice with boundedness again, we find an affine cover $D(a_1),\dots,D(a_n) \subseteq X$ together with sections $h_i : \Gamma(D(a_i), \mc F) \otimes_{R^{D(a_i)}} \Gamma(D(a_i),\mc G)$ such that $s|_{D(a_i)} = \vartheta_{D(a_i)}(h_i)$.
  As above, let us identify $H[a^{-1}] \simeq \Gamma(D(a),\mc F) \otimes_{A[a^{-1}]} \Gamma(D(a), \mc G)$ for all $a:A$.
  Then we regard $h_i$ as an element of $H[a_i^{-1}]$ and  the map $\vartheta_{D(a_ia_j)}$ is of type $H[(a_ia_j)^{-1}] \to \Gamma(D(a_ia_j),\mc F \otimes \mc G)$.
  We have just verified that it is injective, thus the images of $h_i$ and $h_j$ in $H[(a_ia_j)^{-1}]$ agree and glue to an element $h : H$ with $\vartheta_X(h)|_{D(a_i)} = \vartheta_{D(a_i)}(h_i) = s|_{D(a_i)}$ for all $i$.
  Therefore also $\vartheta_X(h) = s$.
\end{proof}

\begin{corollary}
  We have the following.
  \begin{enumerate}
  \item[i)]
    \label{wqc-closed-under-tensorprod}
    If $M$ and $N$ are (weakly) quasi-coherent $R$-modules, then so is $M \otimes N$.

  \item[ii)]
    \label{wqc-closed-under-dep-prod}
    If $\mc F : X \to \Mod R$ is a bundle of (weakly) quasi-coherent $R$-modules on an affine scheme $X$, then $\Gamma(X,\mc F)$ is (weakly) quasi-coherent.
  \end{enumerate}
\end{corollary}

\begin{proof}
  The first item is clear.
  For the second one let $B$ be an arbitrary finitely presented $R$-algebra and suppose $X = \Spec A$.
  Note that $B$ is quasi-coherent as an $R$-module.
  In the following we conflate notation for $B$ and its induced constant bundle $x \mapsto B$ on $X$.
  Consider the composition of isomorphisms:
  \begin{align*}
    \Gamma(X, \mc F) \otimes_R B &\overset{\cong}{\rightarrow} \Gamma(X,\mc F) \otimes_A (A \otimes_R B) \\
                                 &\overset{\cong}{\rightarrow} \Gamma(X,\mc F) \otimes_A \Gamma(X,B)
                                 &&\text{by quasi-coherence of $B$} \\
                                 &\overset{\cong}{\rightarrow} \Gamma(X, \mc F \otimes B)
                                 &&\text{by \Cref{tensor-gamma-commute-on-affines}} \\
                                 &\overset{\cong}{\rightarrow} \Gamma(X, {\mc F}^{\Spec B})
                                 &&\text{by quasi-coherence of $\mc F_x$} \\
                                 &\overset{\cong}{\rightarrow} \Gamma(X,\mc F)^{\Spec B}
  \end{align*}
  We omit the verification that this composition is the given map in the definition of (weak) quasi-coherence.
  This shows the claim for quasi-coherence and, by restricting $B$ to algebras of the form $R[r^{-1}]$, also for weak quasi-coherence.
\end{proof}

\rednote{TODO: Write a warning that \(f_\ast\) and \(f^\ast\) do not have the same properties as in traditional algebraic geometry.}

\begin{lemma}[Projection Formula I]
  \label{projection-formula-I}
  Let $f : X \to Y$ be an affine map of schemes and suppose $\mc F : X \to \Mod R$ is a weakly quasi-coherent and $\mc G : Y \to \Mod R$ a quasi-coherent bundle.
  Then there is an isomorphism of bundles
  \[
    (f_\ast\mc F) \otimes \mc G
    \overset{\cong}{\longrightarrow}
    f_\ast(\mc F \otimes f^\ast \mc G)
  \]
  that is natural in both $\mc F$ and $\mc G$.
\end{lemma}

\begin{proof}
  We verify the statement pointwise for $y:Y$.
  Then consider the following chain of isomorphisms:
  \begingroup
  \let\displaystyle\textstyle
  \begin{align*}
    (f_\ast\mc F \otimes \mc G)_y &= ( \prod_{x:\fib_f(y)} \mc F_x) \otimes \mc G_y \\
                                  &\overset{\cong}{\rightarrow} ( \prod_{x:\fib_f(y)} \mc F_x) \otimes_{R^{\fib_f(y)}} (R^{\fib_f(y)} \otimes \mc G_y) \\
                                  &\overset{\cong}{\rightarrow} ( \prod_{x:\fib_f(y)} \mc F_x) \otimes_{R^{\fib_f(y)}} (\prod_{x:\fib_f(y)} \mc G_y)
                                  && \text{by quasi-coherence of $\mc G_y$} \\
                                  &\overset{\cong}{\rightarrow} \prod_{x:\fib_f(y)} (\mc F_x \otimes \mc G_y)
                                  && \text{by \Cref{tensor-gamma-commute-on-affines}}\\
                                  &= f_\ast (\mc F \otimes f^\ast\mc G)_y
  \end{align*}
  \endgroup
  We omit the calulation showing the naturality claim.
\end{proof}

Let us also prove a second version of the projection formula.
We trade off having to impose a much stronger restriction on $\mc G$ with the benefit of not needing any conditions on $f$ or $\mc F$.

\begin{lemma}[Projection Formula II]
  \label{projection-formula-II}
  Let $f : X \to Y$ be a map of arbitrary types and suppose $\mc F : X \to \Mod R$ and $\mc G : Y \to \Mod R$ are bundles.
  Assume further, that $\mc G$ is a vector bundle.
  Then there is an isomorphism of bundles
  \[
    (f_\ast\mc F) \otimes \mc G
    \overset{\cong}{\longrightarrow}
    f_\ast(\mc F \otimes f^\ast \mc G)
  \]
  that is natural in both $\mc F$ and $\mc G$.
\end{lemma}

\begin{proof}
  Just like the first projection formula we can check this pointwise,
  but using different arguments to show the neccessary identities.
  We omit this standard computation.
\end{proof}

\begin{lemma}
  \label{SES-induced-by-regular-section}
  Suppose $X$ is a scheme.
  Let $\mc L$ be a line bundle and $\mc E$ be a vector bundle on $X$.
  If $s : H^0(X,\mc L)$ is a regular section, then the sequence
  \[
    0 \longrightarrow \mc E \otimes \mc L^\ast
      \overset{\cdot s}{\longrightarrow}  \mc E
      \longrightarrow \mc E^{s=0}
      \longrightarrow 0
  \]
  is affine-locally exact.
  Here, $\mc E^{s=0}$ is the bundle $x \mapsto \mc E_x^{s(x)=0}$.
\end{lemma}

\begin{proof}
  Let $U = \Spec A \subseteq X$ be an affine open subscheme that trivializes both $\mc E$ and $\mc L$, such that $s|_U: \Gamma(U,\mc L) \simeq A$ is a regular element.
  Then, under these trivializations, the sequence $0 \to \Gamma(U,\mc E \otimes \mc L^\ast) \to \Gamma(U, \mc E) \to \Gamma(U,\mc E^{s=0}) \to 0$ is isomorphic to the sequence
  \[
    0 \longrightarrow A^{\rk(\mc E)}
    \overset{s|_U}{\longrightarrow} A^{\rk(\mc E)}
    \longrightarrow \big( A/(s|_U) \big)^{\rk(\mc E)}
      \longrightarrow 0 \rlap{,}
  \]
  which is exact by regularity of $s$.
\end{proof}

\begin{remark}
  \Cref{SES-induced-by-regular-section} can be generalized beyond vector bundles using the notion of $s$-torsion.
  We won't need such a generalization, however.
\end{remark}





\subsection{Basics on Projective Space}

We use notation from \cite{draft} and \cite{sag-projective}.

\begin{definition}
  \label{Pn-definitions}
  $\bP^n$ may be defined as one of the following equivalent types:
  \begin{enumerate}[(i)]
  \item The type of lines\footnote{Submodules $M\subseteq R^{n+1}$ with $\| M=R^1 \|$} through the origin in $R^{n+1}$.
  \item The set-quotient of $R^{n+1}\setminus\{0\}$ by the relation $x\sim y$ iff $x=\lambda y$ for some $\lambda :R^\times$.
  \item\label{Pn-as-homotopy-quotient} The homotopy quotient:
    \[
      \sum_{l : K(R^\times,1)} l^{n+1}\setminus\{0\}
    \]
  \end{enumerate}
\end{definition}

\begin{lemma}
  \label{D-for-Pn}
  Let $U\subseteq \bP^n$.
  The following are equivalent:
  \begin{enumerate}[(i)]
  \item $U$ is open.
  \item $U=D(P_0,\dots,P_l)$ for $P_i:R[X_0,\dots,X_n]_d$.
  \item There are $N:\N$ and $s_0,\dots,s_l:\prod \mathcal O(1)^{\otimes N}$ such that\index{$D(s_0,\dots,s_l)$}
    \[
      U=D(s_0,\dots,s_l)\colonequiv\sum_{x:X}s_0(x)\neq 0\vee \dots \vee s_l(x)\neq 0
      \rlap{.}
    \]
  \end{enumerate}
\end{lemma}

The equivalence between the last two statements follows from \cite{cech-draft}[Theorem 8.3].
The following is a direct consequence of \Cref{Pn-definitions} \ref{Pn-as-homotopy-quotient}:

\begin{remark}
  The type of maps $X\to \bP^n$ is equivalent to the type
  \[
  (\mathcal L:X\to K(R^\times,1))\times (s_0,\dots,s_n:(x:X)\to \mathcal L_x) \to (x:X) \to \exists_{i} s_i(x)\neq 0
  \]
  i.e.\ the type of line bundles on $X$ that come with $n+1$ sections that do not simultaneously vanish.
\end{remark}

\begin{definition}
  A line bundle $\mathcal L:X\to K(R^\times,1)$ is called \notion{very ample}
  if one of the following equivalent statement holds:
  \begin{enumerate}[(i)]
  \item There is a closed embedding $i:X\to \bP^n$ and $\mathcal L=i_\ast \mathcal O(1)$.
  \item There are non-simultaneously vanishing sections $s_0,\dots,s_n$ of $\mathcal L$
    such that $[s_0,\dots,s_n]:X\to \bP^n$ is a closed embedding.
  \end{enumerate}
\end{definition}

\begin{proof}
  All we have to do is to show that if (ii) holds, for $i\colonequiv [s_0,\dots,s_n]$,
  there is an equality $\mathcal L=i_\ast \mathcal O(1)$.
  Let $x:X$ and note first that $\mathcal O(-1)_x$ is generated by $(s_0,\dots,s_n)^T$.
  Let $\varphi:R^1\to \mathcal L_x$ be an isomorphism of $R$-modules.
  Then there are $\lambda_0,\dots,\lambda_n$ such that $\sum \lambda_i s_i(x) =\varphi(1)$,
  which lets us define:
  \begin{center}
    \begin{tikzcd}
      \mathcal L_x \ar[r] & \mathcal O(1)=\Hom_R(\mathcal O(-1),R) \\
      \lambda \cdot \varphi(1)\ar[mapsto,r] & \sum \lambda_i \varphi^{-1}(s_i(x))
    \end{tikzcd}
  \end{center}
  which is an isomorphism independent of the choice of $\varphi$.
\end{proof}


\rednote{The following definition of ample line bundles is never used.}
This is a variant of \cite{vakil}[Theorem 17.6.2] (which might hold):

\begin{definition}
  Let $X$ be a proper scheme.
  A line bundle $\mathcal L : X\to K(R^\times,1)$ is \notion{ample}
  if one of the following equivalent statements holds:
  \begin{enumerate}[(i)]
  \item There is $N>0$ such that $\mathcal L^{\otimes N}$ is very ample.
  \item Any open subset $U\subseteq X$ is of the form $U=D(s_0,\dots,s_n)$,
    for some $s_i:\mathcal L^{\otimes l_i}(X)$ with $l_i>0$.
  \item Any open subset $U\subseteq X$ is of the form $U=D(s_0,\dots,s_n)$,
    for some $N>0$ and $s_i:\mathcal L^{\otimes N}(X)$.
  \item For any finite type, weakly quasi-coherent module bundle $\mathcal F:X\to \Mod{R}_{\mathrm{wqc,fg}}$,
    there is $N>0$, $j:\N$ and a (pointwise) surjection
    $\mathcal O^{\oplus j}\to \mathcal F\otimes \mathcal L^{\otimes N}$.
  \end{enumerate}
\end{definition}

\begin{definition}
  For a scheme $X$, we call
  \[
    T_X : X \to \Mod{R},\; x \mapsto T_xX
  \]
  the \emph{tangent bundle} and
  \[
    \Omega_X : X \to \Mod{R},\; x \mapsto T_xX^\vee
  \]
  then \emph{cotangent bundle}.
  If $X$ is smooth of dimension $n$, we write $\omega_X :\equiv \det(\Omega_X) :\equiv \bigwedge^n\Omega_X$ for the  \emph{canonical line bundle}.
\end{definition}

\begin{lemma}
  There are pointwise exact sequences of $R$-modules
  \[
    0 \longrightarrow
    \mathcal O_{\bP^n} \longrightarrow
    \mathcal O(1)^{\oplus n+1} \longrightarrow
    T_{\bP^n} \longrightarrow
    0
  \]
  and
  \[
    0 \longrightarrow
    \Omega_{\bP^n} \longrightarrow
    \mathcal O(-1)^{\oplus n+1} \longrightarrow
    \mathcal O_{\bP^n} \longrightarrow
    0,
  \]
  where the maps are to be specified in the proof below.
\end{lemma}

  We give a coordinate free proof of this statement.

\begin{proof}
  {\color{red} TODO: Replace with the more conceptual proof.}
  Let $x :\equiv (l,p) : \bP^n$ be arbitrary.
  The exactness of the second sequence follows from the exactness of the first by dualizing,
  which is exact in this case as $T_x\bP^n$ is a free $R$-module.
  For the first sequence, we represent projective space as the homotopy quotient
  \[
    \bP^n = \sum_{l : K(R^\times,1)} l^{n+1}\setminus\{0\}
  \]
  and define auxiliary maps
  \[
    \pi : l^{n+1}\setminus\{0\} \to \bP^n, q \mapsto (l,q)
  \]
  as well as
  \[
    \alpha : l^{\oplus n+1} \to T_p(l^{n+1}\setminus\{0\}), v \mapsto (r \mapsto p + rv).
  \]
  Note that $\alpha$ is an isomorphism of $R$-modules.
  It suffices to show that the composition
  \[
    l^{\oplus n+1} \overset{\alpha}{\longrightarrow}
    T_p(l^{n+1}\setminus\{0\}) \overset{d_p\pi}{\longrightarrow}
  T_{(l,p)}\bP^n
  \]
  is surjective with kernel generated by $p$, for then $T_x\bP^n = l^{\oplus n+1}/(p)$ and we obtain an exact sequence
  \[
    0 \longrightarrow
    R \overset{1 \mapsto p}{\longrightarrow}
    l^{\oplus n+1} \longrightarrow
    T_x\bP^n \longrightarrow 0.
  \]
  \begin{enumerate}
    \item[i)]
    Note that $\pi$ is smooth, as for any $y : \bP^n$ we have that $\propTrunc{\fib_\pi(y) = R^\times}$ and $R^\times$ is smooth.
    It then follows from \cite[Corollary 4.1.1]{diffgeo-article} that $d_p\pi$ is surjective.


    \item[ii)]
    That the kernel is spanned by $p : l^{\oplus n+1}$ can be seen by virtue of a direct calculation.
    Alternatively, consider the fiber sequence $(\fib_\pi(x),p) \to (l^{n+1}\setminus\{0\},p) \to (\bP^n,x)$.
    It induces an exact sequence
    \[
      0 \longrightarrow
      T_p\fib_\pi(x) \longrightarrow
      T_p(l^{n+1}\setminus\{0\}) \longrightarrow
      T_x\bP^n \longrightarrow
      0,
    \]
    where exactness on the right has been verified in the step before.
    Hence we see that $\ker(d_p\pi \circ \alpha) = T_p\fib_\pi(x)$ is free of rank $1$ and must therefore be generated by the non-zero element $p$.
\end{enumerate}
\end{proof}

\begin{corollary}\label{canonical-bundle-as-twist}
  We have an isomorphism $\omega_{\bP^n} = \mathcal O(-n-1)$ of $R$-module bundles.
\end{corollary}

\begin{proof}
  By virtue of the previous lemma we have
  \[
    \omega_{\bP^n} = \det(\Omega_X) \otimes \det(\mathcal O_{\bP^n}) = \det(\mathcal O(-1)^{\oplus n+1}) = \mathcal O (-n-1).
  \]
\end{proof}





\subsection{Étale-Overt Types}

In the following we use \((\_)_{\et}\) to denote étale sheafification, which is the localization at those types $\|X\|$, where $X$ is an affine, formally étale and non-empty scheme.
In particular, $\|X\|_\et$ is the étale sheafification of the propositional truncation of $X$.

Étale sheaves in synthetic algebraic geometry have been extensively studied by Hugo Moeneclay and Thierry Coquand and we build the discussion in this section on their foundational work {\color{red}CITE}.

\begin{definition}
  A type \(X\) is called \emph{étale-overt} if for any open subset \(U\subseteq X\)
  the proposition \(\|U\|_\et\) is open.
\end{definition}

\begin{example}
  \begin{enumerate}
    \item[i)] All open propositions are étale-overt.
    \item[ii)] It is not the case that all closed propositions are étale-overt.
      Indeed, assume the family of propositions $r=0$ is étale-overt for all $r:R$.
      Then $r=0$ would also be an open proposition and we thus would have $\neg\neg(r=0) \to (r=0)$, showing that any nilpotent element of $R$ is zero. This violates duality.
  \end{enumerate}
\end{example}

\begin{lemma}
  \label{etale-overtness-is-local}
  Let $X$ be a type with an open cover $U_1,\dots,U_n$.
  Then $X$ is étale-overt if and only if all $U_i$ are étale-overt.
\end{lemma}

\begin{proof}
  This follows immediately from the equivalence
  \[
    \|X\|_\et = \|U_1\|_\et \lor \dots \lor \|U_n\|_\et\rlap{.}
  \]
\end{proof}

\begin{proposition}
  \label{étale-overt-is-sigma-closed}
  Étale-overt types are closed under dependent sums.
\end{proposition}

\begin{proof}
  Let $Y_x$ be a family of étale-overt types indexed by another étale-overt type $X$.
  Suppose $U \subseteq \sum_{x:X} Y$ is open.
  Writing $U = \sum_{x:X}U_x$ for open subsets $U_x \subseteq Y_x$ we obtain
  \[
    \|U\|_\et = \Big\|\sum_{x:X} \|U_x\|_\et\Big\|_\et\rlap{.}
  \]
  Indeed, the universal properties of propositional truncation and étale-sheafification immediately give a map from left to right. The same argument twice, together with currying the dependent sum after the first application, gives a map from right to left.
  Since $\sum_{x:X} \|U_x\|_\et$ is an open subtype of the étale-overt type $X$, we see that $\|U\|_\et$ is open.
\end{proof}

We will use the following lemma to show that $R^n$ and $\bP^n$ are étale-overt.
A much less general lemma would suffice for this objective alone, this stronger version will however later be instrumental in showing that a smooth projective scheme $X$ admits a finite map $X \to \bP^{\dim X}$; see \Cref{merely-finite-map}.

\begin{lemma}
    \label{consume-non-emptiness-of-open-subset}
  Let $U \subseteq \bP^n$ be open such that $\neg\neg U$.
  Then there merely exists an étale $R$-algebra $A$, that is finite free of positive rank as an $R$-module, together with a map
  \[
    \Spec(A) \longrightarrow U \rlap{.}
  \]
\end{lemma}

\begin{proof}
  Writing $U = D_+(F_1, \hdots, F_m)$, we see that $\neg\neg U$ is equivalent to $\bigvee_j F_j \neq 0$.
  In particular, $\neg\neg U$ is equivalent to $\bigvee_j \neg\neg D_+(F_j)$.
  Thus we may reduce to the case of $m=1$.
  Argueing similarly, we may even replace $\bP^n$ by $R^n$.
  Hence we may assume $U = D(f)$ for some primitive polynomial $f : R[X_1, \hdots, X_n]$.
  Now choose $M$ such that $M > \deg f$ and define a polynomial $F : R[T]$ via $F(T) : \equiv f(T, T^M, \hdots, T^{M^{n-1}})$.
  By choice of $M$, the set of coefficients appearing in $F$ agrees with the set of coefficients of $f$, thus $F$ is primitive as well.
  Suppose $D \geq \deg F$ and let $p:\bN$ be a prime such that $p \geq D + 2$ and $p \neq 0$ in $R$.
  Now we claim that
  \[
    A :\equiv R[X]/(\Phi_p(X)),
    \quad \text{with} \quad
    \Phi_p(X) :\equiv 1 + X + \hdots + X^{p-1} : R[X]
  \]
  satisfies the statement.
  Clearly $A$ is finite free of rank $p - 1$ and a routine verification yields that $\Spec(A)$ is standard étale and hence, of course, étale.
  It remains to construct a map $\Spec(A) \to U$.
  For this we write $x$ for the image of $X$ in $A$.
  Note that $x^k - x^l = x^k(1-x^{l-k})$ is invertible in $A$ for $1 \leq k < l < p$.
  Indeed, since $(X-1)\Phi_p(X) = X^p-1$, we have $x^p = 1 : A$ and thus $x^k$ is invertible in $A$ as well.
  It is now enough to show that $1 - x^l$ is invertible for $1 \leq l < p$.
  As
  \[
    p - \Phi_p(X) = (1 - X) \sum_{j=1}^{p-1} 1 + X + \hdots + X^{j-1} : R[X] \rlap{,}
  \]
  with $\Phi_p(x) = 0 : A$ and $p : A^\ast$, $1-x$ is invertible in $A$.
  Then $1-x^l$, being the image of $1-x$ along the well-defined $R$-algebra map $A \to A, x \mapsto x^l$, is invertible as well.
  Now consider $F$ as a polynomial in $A[T] \equiv (R[X]/(\Phi_p(X))[T]$ and write $F(T) = \sum_{j=0}^{D} c_jT^j$.
  Then the Vandermonde matrix in the linear system
  \[
    \begin{pmatrix}
      1 & 1 & 1 & \cdots & 1 \\
      1 & x & x^2 & \cdots & x^D \\
      \vdots & \vdots & \vdots & \ddots & \vdots\\
      1 & x^D & x^{2D} & \cdots & x^{D^2}
      \end{pmatrix}
      \begin{pmatrix}
      c_0\\
      c_1\\
      \vdots\\
      c_D
    \end{pmatrix}
    =
    \begin{pmatrix}
      F(1) \\
      F(x) \\
      \vdots\\
      F(x^D)
    \end{pmatrix}
    : A^{D+1}
  \]
  is invertible, meaning that the coefficients of $F$ can be expressed as $A$-linear combinations of the $F(x^j)$.
  Thus $(F(1), \hdots, F(x^D)) = A$.
  The ring $A$, being a finite extension of the local ring $R$, is a local-global ring \cite[Chapter IX, Theorem 6.13]{lombardi-quitte}.
  Hence there merely is an $a:A$ such that $F(a)$ is a unit of $A$.
  Define an $R$-algebra map $\alpha : R[X_1, \hdots, X_n] \to A, X_j \mapsto a^{M^{j-1}}$.
  Then $\alpha(f) = F(a)$, hence we get an induced map in
  \[
    \begin{tikzcd}
      R[X_1, \hdots, X_n] \ar[r, "\alpha"] \ar[d] & A \\
      R[X_1, \hdots, X_n][f^{-1}] \ar[ur, "\exists !\overline{\alpha}"', dashed] &
    \end{tikzcd}
  \]
  as desired.

\end{proof}

\begin{corollary}
  \label{R^n-etale-overt}
  Both $\bP^n$ and $R^n$ are étale-overt.
\end{corollary}

\begin{proof}
  Since open subsets of étale-overt types are again étale-overt, it suffices to show that $\bP^n$ is étale-overt.
  Say $U \subseteq \bP^n$ is open.
  It suffices to show that $\|U\|_\et = \neg\neg U$, since $\neg\neg U$ is open as observed in the beginning of the proof of \Cref{consume-non-emptiness-of-open-subset}.
  This fact also makes the implication $\|U\|_\et \to \neg\neg U$ immediate, since open propositions are étale sheaves.
  To verify the converse, assume $\neg\neg U$. Then by \Cref{consume-non-emptiness-of-open-subset} we find a map $\Spec(A) \to U$ and thus in particular maps $\|\Spec(A)\| \to \|U\| \to \|U\|_\et$.
  Since $\Phi_p(X)$ has a zero up to double-negation, we have $\neg\neg\Spec(A)$.
  Hence $\|\Spec(A)\|$ is part of the étale topology.
  Therefore $\|\Spec(A)\| \to \|U\|_\et$ is equivalent to $\|U\|_\et$ and we win.
\end{proof}

The argument in \Cref{R^n-etale-overt} proving that $\neg\neg U$ is open, is very special to the case of projective and affine space.
As far as the author knows, it does not generalize well.
However, as the following lemma shows, even though the above proof does not generalize, the phenomenon does.
This was first proven in \cite[Lemma 1.1.7]{étale-sheaves-compendium}.
We, nonetheless, present a different interesting proof here using Cayley-Hamilton.

\begin{lemma}
  \label{notnot-Spec(R[X]/P)_G-is-open}
  Let $P,G:R[x]$ such that $P$ is monic and non-constant and set $X \equiv \Spec ((R[x]/P)[G^{-1}])$.
  Then $\neg\neg X$ is an open proposition.
\end{lemma}

\begin{proof}
  Let $M :\equiv R[x]/P$ and note that this is a finite free $R$-module of rank $n \equiv \deg P$.
  Define an $R$-linear endomorphism $f : M \to M , m \mapsto G \cdot m$ and denote by
  \[
    \chi(\lambda) = \lambda^n + c_1\lambda^{n-1} + \dots + c_n : R[\lambda]
  \]
  its characteristic polynomial.
  Cayley-Hamilton establishes that $\chi (f) = 0$ and thus in particular
  \[
    0 = \chi (f(1)) = G^n + c_1G^{n-1} + \dots + c_n : M\rlap{.}
  \]
  For $D :\equiv c_1 \neq 0 \lor \dots \lor c_n \neq 0$ we now claim that $D = \neg\neg X$.
  To see $\neg\neg X \to D$ we can assume that we have $x:X$ such that $P(x) = 0$ and $G(x)\neq 0$ and thus get a well-defined equation
  \[
    0 \neq G(x)^n = -c_1G(x)^{n-1} - \dots - c_n
  \]
  in $R$, showing that some $c_j$ must be invertible.
  For the converse direction $D \to \neg\neg X$ we show $\neg X \to \neg D$.
  If $X$ is empty, then $G$ is nilpotent in $M$ and hence $f : M \to M$ is a nilpotent endomorphism, say $f^{m+1}=0$.
  Then we have $(id - \lambda f)(1 + \lambda f + ... + \lambda^m f^m) = id$ and thus $\det(id - \lambda f) = c_n\lambda^n + \dots c_1 \lambda + 1$ is invertible in $R[\lambda]$.
  Therefore the coefficients $c_1,\dots,c_n$ are nilpotent as well, which is to say that $\neg D$.
\end{proof}

\begin{lemma}
  An étale scheme is étale-overt.
\end{lemma}

\begin{proof}
  Being étale-overt is a local property by \Cref{etale-overtness-is-local}  and étale schemes are locally basic étale, so it is enough to show that any basic étale scheme $X = \Spec((R[X]/P)[G^{-1}])$ is étale-overt.
  Since even any open $U \subseteq X$ is étale again, the same argument shows that it is enough to prove that $\|X\|_\et$ is open.
  From \Cref{notnot-Spec(R[X]/P)_G-is-open} we know that $\neg\neg X$ is open, so we are done once we verify that
  \[
    \neg\neg X = \|X\|_\et\rlap{.}
  \]
  As open propositions are étale sheaves, we immediately get a map from right to left. For the opposite direction note that once we assume $\neg\neg X$, $\|X\|$ becomes part of the étale topology and hence
  $\|X\| \to \|X\|_\et$ implies $\|X\|_\et$, but the former proposition is clearly true.
\end{proof}

\begin{corollary}
  \label{smooth-implies-etale-overt}
  Any smooth scheme is étale-overt.
\end{corollary}

\begin{proof}
  Let $X$ be a smooth scheme.
  Because of \Cref{etale-overtness-is-local} we may assume that $X$ is standard smooth.
  Then we merely have an étale map $f : X \to R^n$, i.e.\ $X = \sum_{x:R^n} \fib_f(x)$.
  Hence $X$, being a dependent sum of étale-overt types, is itself étale-overt by \Cref{étale-overt-is-sigma-closed}.
\end{proof}




\subsection{Finite and Finite Free Maps}

\begin{definition}
  \label{def:finite-free-map}
  Let $f : X \to Y$ be a map of schemes such that all fibers of \(f\) are affine schemes.
  We say that $f$ is \emph{finite}, if $f_\ast\mc O_X$ is a bundle of pointwise finitely presented $R$-modules.
  If $f_\ast\mc O_X$ is even a vector bundle (of degree $d$), then we call $f$ \emph{finite free (of degree $d$)}.
  Furthermore, we say a scheme $X$ is \emph{finite (free)}, if $X \to 1$ is finite (free).
\end{definition}

It is important to distinguish the concept of a finite (free) scheme, as just defined, from the concept of a \emph{standard finite} type.
Recall that this is a type $X$, such that
\[
  \Big\| \sum_{n:\bN} X = \{1,\hdots,n\} \Big\| \rlap{.}
\]
A scheme that is a standard finite type is finite free, and a finite free scheme is finite.
However, no reverse implication holds.

In classical algebraic geometry, a morphism of schemes $f : X \to Y$ such that $f_\ast\mc O_X$ is locally finite free is called \emph{finite locally free}.
We intentionally drop the adverb \emph{locally} from our convention, as our definition is a pointwise, rather than an affine-local, condition.

\begin{lemma}
  Let $f : X \to Y$ be a map of schemes.
  Then $f$ is finite free if and only if it is finite and flat.
\end{lemma}

\begin{proof}
  By \cite[Theorem VIII-1.4]{lombardi-quitte} the $R$-module $(f_\ast \mc O_X)_y = R^{\fib_f(y)}$ is finitely generated and projective if and only if it is finitely presented and flat.
  Since finitely generated projective modules over local rings are precisely the finite free modules, we are done.
\end{proof}

\begin{lemma}
  \label{finite-free-morphisms-preserve-vector-bundles}
  Let $f : X \to Y$ be a finite free morphism between schemes and suppose $\mc E$ is a vector bundle on $X$.
  Then $f_\ast \mc E$ is a vector bundle as well.
\end{lemma}

\begin{proof}
  First note that $f_\ast\mc E$ is an $f_\ast\mc O_X$-algebra and abbreviate $A :\equiv R^{\fib_f(y)}$.
  We know that $(f_\ast\mc E)_y = \prod_{x:\fib_f(y)}\mc E_x$ is a finitely generated projective $A$-module.
  By assumptions on $f$, we have that $A$ is finite free as an $R$-module.
  These two facts together imply that $(f_\ast \mc E)_y$ is even a finitely generated projective $R$-module and hence is finite free over $R$.
\end{proof}

\begin{theorem}
  \label{miracle-flatness}
  Let $R$ be any commutative ring with $1$ and let $A$ be an $R$-algebra that is finite.
  Suppose there merely are $f_1,\dots,f_n:R[x_1,\dots,x_n]$ such that
  \[
    A \cong R[x_1,\dots,x_n]/(f_1,\dots,f_n) \rlap{.}
  \]
  Then $A$ is a projective and flat $R$-module.
\end{theorem}

\begin{proof}
  A constructively valid proof is given in \cite[Théorème 12]{lombardi2025theoremesmitetlenstra}.
\end{proof}


\begin{lemma}
  \label{fin-subscheme-contained-in-std-smooth}
  Let $F$ be a finite and closed subscheme of an affine scheme $X$ that is smooth.
  Then there merely is a function $h : X \to R$ such that $D(h)$ is standard smooth and such that $F \subseteq D(h)$.
\end{lemma}

\begin{proof}
  {\color{red}TODO: Write down my proof.}
\end{proof}

\begin{lemma}
  \label{finite-map-of-equidimensional-smooth-schemes-is-finite-free}
  Let $f : X \to Y = \bP^n$ be a finite map between smooth schemes $X$ and $Y$ of the same dimension $n:\bN$.
  Then $f_\ast\mc O_X$ is a vector bundle.
\end{lemma}

{\color{red}TODO: generalize from $\bP^n$ to an arbitrary smooth $Y$ with $\dim Y = n$.}

\begin{proof}
  Without loss of generality replace $\bP^n$ by $R^n$.
  Since $f$ is affine and affine schemes are closed under dependent sums, we thus see that $X$ is affine.
  It is enough to show that $(f_\ast\mc O_X)_0 = R^{\fib_f(0)}$ is finite free.
  By \Cref{fin-subscheme-contained-in-std-smooth} we find a function $h : X \to R$ such that $D(h)$ is standard smooth of dimension $n$ and $\fib_f(0) \subseteq D(h)$.
  Suppose that $D(h) = \Spec R[X_1,\hdots, X_m, Y_1, \hdots, Y_n]/(P_1, \hdots P_m)$.
  Then
  \[
    R^{\fib_f(0)} \cong R[X_1, \hdots, X_m, Y_1, \hdots, Y_n]/(P_1, \hdots, P_m, f_1, \hdots, f_n) \rlap{.}
  \]
  This $R$-algebra is finite by assumption, hence by \Cref{miracle-flatness} we see that $R^{\fib_f(0)}$ is projective and thus finite free.
\end{proof}

\begin{corollary}
  \label{finite-map-of-equidimensional-smooth-schemes-preserves-vectorbundles}
  In the situation of \Cref{finite-map-of-equidimensional-smooth-schemes-is-finite-free}, the pushforward $f_\ast$ preserves vector bundles.
\end{corollary}

\begin{proof}
  This follows directly from \Cref{finite-free-morphisms-preserve-vector-bundles}.
\end{proof}

